% preamble-exam-zh.tex — 共享 preamble
% 用于所有 exam-zh 模板
%
% 样式策略：
%   屏幕 → 语义彩色（蓝=关键想法，红=易错，绿=路线，灰=提示/教师备注）
%   打印 → 自动降级为黑白（通过 print mode 或 monochrome 选项）

\documentclass{exam-zh}

% ---------- 基础配置 ----------
\examsetup{
  page/size = a4paper,
  page/show-head = false,
  page/show-foot = false,
  sealline/show = false,
  font = times,
  math-font = xits,
}

\usepackage{tcolorbox}
\usepackage{needspace}
\usepackage{graphicx}
\tcbuselibrary{skins,breakable}

% ---------- 打印/屏幕颜色切换 ----------
% 用法：\documentclass[monochrome]{exam-zh} 即可切换为纯黑白
% 注意：不要使用 \if@xxx 放在 makeatother 外部；这里改成普通条件名。
\newif\ifeduprintcolor
\eduprintcolortrue

\makeatletter
\@ifclasswith{exam-zh}{monochrome}{\eduprintcolorfalse}{}
\makeatother

% ---------- 语义颜色定义 ----------
% 屏幕：有色彩区分；打印：降级为灰度
\ifeduprintcolor
  % --- 屏幕彩色模式 ---
  \definecolor{edu-blue}{RGB}{47, 82, 122}       % 关键想法
  \definecolor{edu-blue-bg}{RGB}{243, 246, 252}
  \definecolor{edu-red}{RGB}{180, 40, 40}         % 易错提醒
  \definecolor{edu-red-bg}{RGB}{255, 245, 240}
  \definecolor{edu-green}{RGB}{34, 100, 60}       % 路线图
  \definecolor{edu-green-bg}{RGB}{242, 250, 245}
  \definecolor{edu-orange}{RGB}{160, 90, 0}       % 训练目标
  \definecolor{edu-orange-bg}{RGB}{255, 248, 235}
  \definecolor{edu-gray}{RGB}{100, 100, 100}      % 提示/教师备注
  \definecolor{edu-gray-bg}{RGB}{248, 248, 248}
  \definecolor{edu-purple}{RGB}{90, 50, 120}      % 思考题
  \definecolor{edu-purple-bg}{RGB}{248, 244, 252}
\else
  % --- 打印黑白模式 ---
  \definecolor{edu-blue}{gray}{0.15}
  \definecolor{edu-blue-bg}{gray}{0.96}
  \definecolor{edu-red}{gray}{0.25}
  \definecolor{edu-red-bg}{gray}{0.97}
  \definecolor{edu-green}{gray}{0.20}
  \definecolor{edu-green-bg}{gray}{0.97}
  \definecolor{edu-orange}{gray}{0.25}
  \definecolor{edu-orange-bg}{gray}{0.98}
  \definecolor{edu-gray}{gray}{0.40}
  \definecolor{edu-gray-bg}{gray}{0.97}
  \definecolor{edu-purple}{gray}{0.30}
  \definecolor{edu-purple-bg}{gray}{0.98}
\fi
\colorlet{edu-diagram-muted}{edu-gray}

% ---------- 自定义教学环境 ----------
% 共性：enhanced + breakable（跨页不断裂）

% 原题卡片 — 强边框，突出显示
\newtcolorbox{problemcard}{
  enhanced, breakable,
  colback=edu-blue-bg,
  colframe=edu-blue,
  boxrule=1.2pt,
  arc=0pt,
  left=3mm, right=3mm, top=3mm, bottom=3mm,
}

% 解题路线图 — 左边线 + 浅底
\newtcolorbox{routebox}{
  enhanced, breakable,
  colback=edu-green-bg,
  colframe=edu-green!30,
  boxrule=0pt,
  borderline west={2.5pt}{0pt}{edu-green},
  left=3mm, right=3mm, top=3mm, bottom=3mm,
}

% 关键想法 — 蓝色左边线
\newtcolorbox{keyideabox}{
  enhanced, breakable,
  colback=edu-blue-bg,
  colframe=edu-blue!30,
  boxrule=0pt,
  borderline west={2.5pt}{0pt}{edu-blue},
  left=3mm, right=3mm, top=3mm, bottom=3mm,
}

% 易错提醒 — 红色左边线
\newtcolorbox{mistakebox}{
  enhanced, breakable,
  colback=edu-red-bg,
  colframe=edu-red!20,
  boxrule=0pt,
  borderline west={3pt}{0pt}{edu-red},
  left=3mm, right=3mm, top=2.5mm, bottom=2.5mm,
}

% 提示 — 灰色虚线左边线
\newtcolorbox{hintbox}{
  enhanced, breakable,
  colback=edu-gray-bg,
  colframe=edu-gray!20,
  boxrule=0pt,
  borderline west={2pt}{0pt}{edu-gray, dashed},
  left=3mm, right=3mm, top=2mm, bottom=2mm,
}

% 思考题 — 紫色虚线左边线
\newtcolorbox{thinkbox}{
  enhanced, breakable,
  colback=edu-purple-bg,
  colframe=edu-purple!20,
  boxrule=0pt,
  borderline west={2pt}{0pt}{edu-purple, dashed},
  left=2.5mm, right=2.5mm, top=2.5mm, bottom=2.5mm,
}

% 教师备注 — 灰色左边线，小字
\newtcolorbox{teachernote}{
  enhanced, breakable,
  colback=edu-gray-bg,
  colframe=edu-gray!15,
  boxrule=0pt,
  borderline west={2pt}{0pt}{edu-gray!60},
  left=2.5mm, right=2.5mm, top=2.5mm, bottom=2.5mm,
  fontupper=\small,
  colupper=black!60,
}

% 训练目标 — 橙色左边线，小字
\newtcolorbox{traininggoal}{
  enhanced, breakable,
  colback=edu-orange-bg,
  colframe=edu-orange!15,
  boxrule=0pt,
  borderline west={1.5pt}{0pt}{edu-orange},
  left=2mm, right=2mm, top=1mm, bottom=1mm,
  fontupper=\small,
}

% 答题区 — 无边框，纯留白
\newcommand{\answerarea}[1][40mm]{%
  \vspace{2mm}%
  \begin{tcolorbox}[
    enhanced, breakable,
    colback=white,
    colframe=white,
    boxrule=0pt,
    height=#1,
    valign=top,
    bottom=0mm,
    after skip=0mm,
  ]
  \end{tcolorbox}%
}

\newcommand{\diagramcolfigure}[3]{%
  \begin{minipage}[t]{#2}
    \vspace{0pt}\centering
    \includegraphics[width=\linewidth]{\detokenize{#1}}
    \if\relax\detokenize{#3}\relax\else
      \par\vspace{1mm}{\scriptsize\color{edu-diagram-muted}#3}
    \fi
  \end{minipage}%
}

\newcommand{\diagramrowitem}[4]{%
  \begin{minipage}[t]{#2}
    \vspace{0pt}\centering
    \includegraphics[width=\linewidth]{\detokenize{#1}}
    \if\relax\detokenize{#3}\relax\else
      \par\vspace{1mm}{\scriptsize\bfseries #3}
    \fi
    \if\relax\detokenize{#4}\relax\else
      \par\vspace{0.5mm}{\scriptsize\color{edu-diagram-muted}#4}
    \fi
  \end{minipage}%
}

\newcommand{\answerareawithdiagramcol}[4]{%
  \vspace{2mm}%
  \begin{tcolorbox}[
    enhanced, breakable,
    colback=white,
    colframe=white,
    boxrule=0pt,
    height=#1,
    valign=top,
    bottom=0mm,
    after skip=0mm,
  ]
  \begin{minipage}[t]{\dimexpr\linewidth-#3-6mm\relax}
    \vspace{0pt}
  \end{minipage}\hfill
  \diagramcolfigure{#2}{#3}{#4}
  \end{tcolorbox}%
}

\newcommand{\answerareawithdiagram}[4]{%
  \answerareawithdiagramcol{#1}{#2}{#3}{#4}%
}

% ---------- 字体设置 ----------
% exam-zh (ctexbook) already sets CJK fonts via ctex.
% Only override if specific fonts are needed.
% macOS: Songti SC, STHeiti, STFangsong
% Linux: SimSun, SimHei, FangSong

% ---------- 页面设置 ----------
% exam-zh (ctexbook) already loads geometry.
% Override margins if needed:
% \geometry{top=16mm, bottom=16mm, left=14mm, right=14mm}

\examsetup{
  question/show-answer = true,
  fillin/show-answer = true,
  solution/show-solution = show-stay,
}

% ---------- 教师版解答样式 ----------
\newtcolorbox{solutionblock}{
  enhanced, breakable,
  colback=edu-blue-bg,
  colframe=edu-blue!25,
  boxrule=0pt,
  borderline west={2pt}{0pt}{edu-blue!50},
  arc=0pt,
  left=3mm, right=3mm, top=2mm, bottom=2mm,
  before skip=2mm,
  after skip=3mm,
  fontupper=\small,
}

% 解答步骤编号
\newcounter{solstep}
\newcommand{\solstep}[2]{%
  \stepcounter{solstep}%
  \par\medskip
  \noindent\hangindent=6mm
  {\color{edu-blue}\bfseries\textcircled{\scriptsize\thesolstep}\enspace #1}\enspace #2\par
}

\begin{document}

\section*{一元二次方程 · 20题方法练习（教师版）}

\section*{一、开平方法（每题 5 分，共 25 分）}

\needspace{8\baselineskip}

\begin{problem}[points=5]
  解方程：$(x-3)^2=16$。
  \setcounter{solstep}{0}
  \begin{solutionblock}
    \textbf{答案：}$x=7$ 或 $x=-1$\par\medskip
    \solstep{开平方}{由 $(x-3)^2=16$ 得 $x-3=\pm4$，所以 $x=7$ 或 $x=-1$。}
  \end{solutionblock}
\end{problem}


\needspace{8\baselineskip}

\begin{problem}[points=5]
  解方程：$(2x+1)^2=25$。
  \setcounter{solstep}{0}
  \begin{solutionblock}
    \textbf{答案：}$x=2$ 或 $x=-3$\par\medskip
    \solstep{开平方}{由 $(2x+1)^2=25$ 得 $2x+1=\pm5$，所以 $x=2$ 或 $x=-3$。}
  \end{solutionblock}
\end{problem}


\needspace{8\baselineskip}

\begin{problem}[points=5]
  解方程：$3x^2-27=0$。
  \setcounter{solstep}{0}
  \begin{solutionblock}
    \textbf{答案：}$x=3$ 或 $x=-3$\par\medskip
    \solstep{整理后开方}{$3x^2=27$，得 $x^2=9$，所以 $x=\pm3$。}
  \end{solutionblock}
\end{problem}


\needspace{8\baselineskip}

\begin{problem}[points=5]
  解方程：$(x+4)^2=7$。
  \setcounter{solstep}{0}
  \begin{solutionblock}
    \textbf{答案：}$x=-4+\sqrt7$ 或 $x=-4-\sqrt7$\par\medskip
    \solstep{开平方}{$x+4=\pm\sqrt7$，所以 $x=-4\pm\sqrt7$。}
  \end{solutionblock}
\end{problem}


\needspace{8\baselineskip}

\begin{problem}[points=5]
  解方程：$5(x-2)^2=45$。
  \setcounter{solstep}{0}
  \begin{solutionblock}
    \textbf{答案：}$x=5$ 或 $x=-1$\par\medskip
    \solstep{先除系数}{两边同除以 $5$，得 $(x-2)^2=9$，所以 $x-2=\pm3$，即 $x=5$ 或 $x=-1$。}
  \end{solutionblock}
\end{problem}


\section*{二、十字相乘法（每题 5 分，共 25 分）}

\needspace{8\baselineskip}

\begin{problem}[points=5]
  解方程：$x^2-7x+12=0$。
  \setcounter{solstep}{0}
  \begin{solutionblock}
    \textbf{答案：}$x=3$ 或 $x=4$\par\medskip
    \solstep{因式分解}{$x^2-7x+12=(x-3)(x-4)$，所以 $(x-3)(x-4)=0$，得 $x=3$ 或 $x=4$。}
  \end{solutionblock}
\end{problem}


\needspace{8\baselineskip}

\begin{problem}[points=5]
  解方程：$x^2+x-20=0$。
  \setcounter{solstep}{0}
  \begin{solutionblock}
    \textbf{答案：}$x=4$ 或 $x=-5$\par\medskip
    \solstep{因式分解}{$x^2+x-20=(x+5)(x-4)$，所以 $x=4$ 或 $x=-5$。}
  \end{solutionblock}
\end{problem}


\needspace{8\baselineskip}

\begin{problem}[points=5]
  解方程：$2x^2+7x+3=0$。
  \setcounter{solstep}{0}
  \begin{solutionblock}
    \textbf{答案：}$x=-3$ 或 $x=-\frac12$\par\medskip
    \solstep{因式分解}{$2x^2+7x+3=(2x+1)(x+3)$，所以 $2x+1=0$ 或 $x+3=0$，得 $x=-\frac12$ 或 $x=-3$。}
  \end{solutionblock}
\end{problem}


\needspace{8\baselineskip}

\begin{problem}[points=5]
  解方程：$3x^2-10x+8=0$。
  \setcounter{solstep}{0}
  \begin{solutionblock}
    \textbf{答案：}$x=2$ 或 $x=\frac43$\par\medskip
    \solstep{因式分解}{$3x^2-10x+8=(3x-4)(x-2)$，所以 $x=\frac43$ 或 $x=2$。}
  \end{solutionblock}
\end{problem}


\needspace{8\baselineskip}

\begin{problem}[points=5]
  解方程：$4x^2-12x+9=0$。
  \setcounter{solstep}{0}
  \begin{solutionblock}
    \textbf{答案：}$x_1=x_2=\frac32$\par\medskip
    \solstep{完全平方}{$4x^2-12x+9=(2x-3)^2$，所以 $2x-3=0$，得 $x_1=x_2=\frac32$。}
  \end{solutionblock}
\end{problem}


\section*{三、配方法（每题 5 分，共 25 分）}

\needspace{8\baselineskip}

\begin{problem}[points=5]
  用配方法解方程：$x^2+6x+5=0$。
  \setcounter{solstep}{0}
  \begin{solutionblock}
    \textbf{答案：}$x=-1$ 或 $x=-5$\par\medskip
    \solstep{配方}{$x^2+6x=-5$，两边加 $9$ 得 $(x+3)^2=4$，所以 $x+3=\pm2$，得 $x=-1$ 或 $x=-5$。}
  \end{solutionblock}
\end{problem}


\needspace{8\baselineskip}

\begin{problem}[points=5]
  用配方法解方程：$x^2-4x-1=0$。
  \setcounter{solstep}{0}
  \begin{solutionblock}
    \textbf{答案：}$x=2+\sqrt5$ 或 $x=2-\sqrt5$\par\medskip
    \solstep{配方}{$x^2-4x=1$，两边加 $4$ 得 $(x-2)^2=5$，所以 $x=2\pm\sqrt5$。}
  \end{solutionblock}
\end{problem}


\needspace{8\baselineskip}

\begin{problem}[points=5]
  用配方法解方程：$x^2+5x+2=0$。
  \setcounter{solstep}{0}
  \begin{solutionblock}
    \textbf{答案：}$x=\frac{-5+\sqrt{17}}2$ 或 $x=\frac{-5-\sqrt{17}}2$\par\medskip
    \solstep{配方}{$x^2+5x=-2$，两边加 $\frac{25}{4}$ 得 $(x+\frac52)^2=\frac{17}{4}$，所以 $x=\frac{-5\pm\sqrt{17}}2$。}
  \end{solutionblock}
\end{problem}


\needspace{8\baselineskip}

\begin{problem}[points=5]
  用配方法解方程：$2x^2-8x+3=0$。
  \setcounter{solstep}{0}
  \begin{solutionblock}
    \textbf{答案：}$x=2+\frac{\sqrt{10}}2$ 或 $x=2-\frac{\sqrt{10}}2$\par\medskip
    \solstep{先化首项系数}{两边同除以 $2$ 得 $x^2-4x+\frac32=0$，即 $x^2-4x=-\frac32$。两边加 $4$ 得 $(x-2)^2=\frac52$，所以 $x=2\pm\frac{\sqrt{10}}2$。}
  \end{solutionblock}
\end{problem}


\needspace{8\baselineskip}

\begin{problem}[points=5]
  用配方法解方程：$x^2-3x-4=0$。
  \setcounter{solstep}{0}
  \begin{solutionblock}
    \textbf{答案：}$x=4$ 或 $x=-1$\par\medskip
    \solstep{配方}{$x^2-3x=4$，两边加 $\frac94$ 得 $(x-\frac32)^2=\frac{25}{4}$，所以 $x=\frac32\pm\frac52$，即 $x=4$ 或 $x=-1$。}
  \end{solutionblock}
\end{problem}


\section*{四、求根公式（每题 5 分，共 25 分）}

\needspace{8\baselineskip}

\begin{problem}[points=5]
  用求根公式解方程：$x^2-2x-3=0$。
  \setcounter{solstep}{0}
  \begin{solutionblock}
    \textbf{答案：}$x=3$ 或 $x=-1$\par\medskip
    \solstep{代入公式}{$a=1,b=-2,c=-3$，$\Delta=b^2-4ac=16$，所以 $x=\frac{2\pm4}{2}$，得 $x=3$ 或 $x=-1$。}
  \end{solutionblock}
\end{problem}


\needspace{8\baselineskip}

\begin{problem}[points=5]
  用求根公式解方程：$2x^2-4x+1=0$。
  \setcounter{solstep}{0}
  \begin{solutionblock}
    \textbf{答案：}$x=1+\frac{\sqrt2}{2}$ 或 $x=1-\frac{\sqrt2}{2}$\par\medskip
    \solstep{代入公式}{$a=2,b=-4,c=1$，$\Delta=16-8=8$，所以 $x=\frac{4\pm\sqrt8}{4}=1\pm\frac{\sqrt2}{2}$。}
  \end{solutionblock}
\end{problem}


\needspace{8\baselineskip}

\begin{problem}[points=5]
  用求根公式解方程：$3x^2+2x-1=0$。
  \setcounter{solstep}{0}
  \begin{solutionblock}
    \textbf{答案：}$x=\frac13$ 或 $x=-1$\par\medskip
    \solstep{代入公式}{$a=3,b=2,c=-1$，$\Delta=4+12=16$，所以 $x=\frac{-2\pm4}{6}$，得 $x=\frac13$ 或 $x=-1$。}
  \end{solutionblock}
\end{problem}


\needspace{8\baselineskip}

\begin{problem}[points=5]
  用求根公式解方程：$x^2+4x+4=0$。
  \setcounter{solstep}{0}
  \begin{solutionblock}
    \textbf{答案：}$x_1=x_2=-2$\par\medskip
    \solstep{判别式为零}{$a=1,b=4,c=4$，$\Delta=16-16=0$，所以 $x=\frac{-4}{2}=-2$，即 $x_1=x_2=-2$。}
  \end{solutionblock}
\end{problem}


\needspace{8\baselineskip}

\begin{problem}[points=5]
  用求根公式解方程：$5x^2+x-2=0$。
  \setcounter{solstep}{0}
  \begin{solutionblock}
    \textbf{答案：}$x=\frac{-1+\sqrt{41}}{10}$ 或 $x=\frac{-1-\sqrt{41}}{10}$\par\medskip
    \solstep{代入公式}{$a=5,b=1,c=-2$，$\Delta=1+40=41$，所以 $x=\frac{-1\pm\sqrt{41}}{10}$。}
  \end{solutionblock}
\end{problem}


\clearpage
\section*{答案速查}


\begin{keyideabox}
  \textbf{20题答案}
  1. $7,-1$；2. $2,-3$；3. $\pm3$；4. $-4\pm\sqrt7$；5. $5,-1$。

6. $3,4$；7. $4,-5$；8. $-3,-\frac12$；9. $2,\frac43$；10. $x_1=x_2=\frac32$。

11. $-1,-5$；12. $2\pm\sqrt5$；13. $\frac{-5\pm\sqrt{17}}2$；14. $2\pm\frac{\sqrt{10}}2$；15. $4,-1$。

16. $3,-1$；17. $1\pm\frac{\sqrt2}{2}$；18. $\frac13,-1$；19. $x_1=x_2=-2$；20. $\frac{-1\pm\sqrt{41}}{10}$。

\end{keyideabox}



\end{document}
